\scanpage{6}
\textit{C. R. Acad. Sc.}, t. \textbf{256}, p. 842-844.\hfill Séance du 21 janvier 1963.\hfill Groupe 1\par\noindent\rule{\linewidth}{0.4pt}

\textsc{Algèbre.} --- \textit{Sur les modules de différentielles.} Note (\textsuperscript{*}) de M. \textbf{\textsc{Jean-Pierre Jouanolou}}, présentée par M. René Garnier.

Application des modules de différentielles aux algèbres quasi-séparables et aux anneaux locaux réguliers.

$\mathrm A$ et $\mathrm B$ désignent dans toute cette Note deux anneaux commutatifs et unitaires, $u:\mathrm A\to\mathrm B$ un monomorphisme unitaire et $\mathrm D_{\mathrm A}(\mathrm B)$ le $\mathrm B$-module des différentielles de l’algèbre ainsi définie. Le passage de cette situation « globale » à la situation « locale » se fait en imposant à $\mathrm A$, $\mathrm B$, $u$ d’être locaux; l’idéal maximal de $\mathrm A$ (resp. $\mathrm B$) sera en ce cas noté $\mathfrak m$ (resp. $\mathfrak n$) et son corps résiduel $\mathrm L_a$ (resp. $\mathrm L_b$); par ailleurs, le rang de $\mathfrak m$ est $\operatorname{rg}\mathfrak m=\dim_{\mathrm L_a}(\mathfrak m/\mathfrak m^2)$, et de même pour $\mathfrak n$. Rappelons qu’étant donnée une extension de corps : $\mathrm K\to\mathrm L$, la dimension du $\mathrm L$-espace vectoriel $\mathrm D_{\mathrm K}(\mathrm L)$ représente en caractéristique $0$ (resp. $p>0$) le degré de transcendance de $\mathrm L$ par rapport à $\mathrm K$ (resp. le nombre d’éléments d’une $p$-base de $\mathrm L$ par rapport à $\mathrm K$) : on la notera $\{\mathrm L:\mathrm K\}$ dans tous les cas.

Nous nous proposons d’étudier d’abord les cas où $\mathrm D_{\mathrm A}(\mathrm B)$ est nul, puis ceux où il est libre (en liaison avec la régularité des anneaux locaux), le tout grâce aux deux suites exactes établies dans le cadre local par Berger et Kunz (\textsuperscript{1}) :
\[
\begin{aligned}
0\to\frac{\mathfrak n}{\mathfrak n^2+\mathfrak m\mathrm B}&\to\mathrm L_b\otimes_{\mathrm B}\mathrm D_{\mathrm A}(\mathrm B)\to\mathrm D_{\mathrm L_a}(\mathrm L_b)\to0,\\
0\to\frac{\mathfrak n}{\mathfrak n^2+(\mathfrak n\cap\mathrm A[\mathrm B^p])\,\mathrm B}&\to\mathrm L_b\otimes_{\mathrm B}\mathrm D_{\mathrm A}(\mathrm B)\to\mathrm D_{\mathrm L_a}(\mathrm L_b)\to0.
\end{aligned}
\]

La première est vraie si $\mathrm L_b$ est séparable sur $\mathrm L_a$, la deuxième, sinon : alors $p>0$ désigne la caractéristique des corps résiduels.

1. \textsc{Structure des algèbres quasi-séparables.} --- Une $\mathrm A$-algèbre $\mathrm B$ est dite \uline{quasi-séparable} lorsque $\mathrm D_{\mathrm A}(\mathrm B)=0$. Il s’agit là d’une propriété locale, en ce sens que $\mathrm B$ est quasi-séparable si pour tout idéal maximal $\mathfrak n$ de $\mathrm B$, $\mathrm B_{\mathfrak n}$ est quasi-séparable sur $\mathrm A_{\mathfrak m}$ ($\mathfrak m=\mathfrak n\cap\mathrm A$), et seulement en ce cas.
% slanted
\margin{\handwritten{\uline{non ramifiée}}}

\textit{Localement :}

(i) Si $\mathrm L_b$ est séparable sur $\mathrm L_a$ (en particulier en caractéristique $0$) : lorsque $\mathrm B$ est quasi-séparable,
\[
\mathfrak n=\mathfrak n^2+\mathfrak m\mathrm B,\tag*{(R)}
\]
et $\mathrm L_b$ est algébrique (resp. relativement parfait (\textsuperscript{2})) sur $\mathrm L_a$ en caractéristique $\mathrm O$ (resp. $p>0$) la réciproque étant vraie si l’on suppose, en outre, que $\mathrm D_{\mathrm A}(\mathrm B)$ est de type fini. Notons que si $\mathrm B$ est nœthérien, ou plus généralement si $\mathfrak n$ est de type fini, la relation (R) s’écrit
\[
\mathfrak n=\mathfrak m\mathrm B,
\]
ce qui exprime que $\mathrm B$ n’est pas ramifié sur $\mathrm A$.
\margin{\handwritten{!?}}
% archive sheet number: 5
